\documentclass[usenames,dvipsnames]{beamer}

\mode<presentation> {
\usetheme{Montpellier}
}

\usepackage{graphicx} % Allows including images
\usepackage{booktabs} % Allows the use of \toprule, \midrule and \bottomrule in tables
\usepackage[frenchb]{babel}
\usepackage[T1]{fontenc}
\usepackage[utf8]{inputenc}
\usepackage{algorithm2e}
\setcounter{tocdepth}{2}
\graphicspath{{../Figures/}}


\DeclareMathOperator*{\argmin}{arg\,min}
\DeclareMathOperator*{\argmax}{arg\,max}
\newtheorem*{proposition}{Proposition}
\newcommand\Fontvi{\fontsize{8}{8}\selectfont}

% exercices comptés
\newcounter{numexos}%Création d'un compteur qui s'appelle numexos
\setcounter{numexos}{0}%initialisation du compteur
\newcommand{\exercice}[1]{%Création d'une macro ayant un paramètre
\addtocounter{numexos}{1}%chaque fois que cette macro est appelée, elle ajoute 1 au compteur numexos
\textcolor{DarkOrchid}{Exercice\,\thenumexos\,:}\,#1%Met en rouge Exercice et la valeur du compteur appelée par \thenumeexos
}

\setbeamertemplate{title page}{
    \centering

    \huge\bfseries
    Fondamentaux théoriques du machine learning


    }

\title[FTML]{title}


\date{} % Date, can be changed to a custom date

\makeatletter
\let\@@magyar@captionfix\relax
\makeatother

\begin{document}

\begin{frame}
\titlepage % Print the title page as the first slide
\end{frame}

\section{Risks and risk decompositions}
\frame{\tableofcontents[currentsection]}

\begin{frame}{Deterministic bound on the estimation error}
We consider the best estimator in the hypothesis space $F$.
    \begin{equation*}
	f_a = \argmin_{h\in F}R(h)
    \end{equation*}

    \exercice{}
    Let us show that
    \begin{equation}
	R( f_n)-R(f_a) \leq 2 \sup_{h\in F}|R(h)-R_n(h)|
    \end{equation}
\end{frame}

\begin{frame}{Deterministic bound on the estimation error}
    \begin{equation*}
	f_a = \argmin_{h\in F}R(h)
    \end{equation*}

    \begin{equation}
        \begin{aligned}
	    R(f_n)-R(f_a) &= \big(R(f_n)-R_n( f_n)\big)\\
	    &+\big(R_n(f_n)- R_n(f_a)\big)\\
	    &+\big(R_n(f_a)-R(f_a)\big)\\
	    &\leq | R(f_n)- R_n(f_n) |\\
	    &+\big(R_n(f_n)-R_n(f_a)\big)\\
	    &+ | R_n(f_a)-R(f_a) | \\
	    &\leq 2 \sup_{h\in F}|R(h)-R_n(h)|\\
	    &+\big(R_n(f_n)-R_n(f_a)\big)
        \end{aligned}
    \end{equation}

    But by definition $ f_n$ minimizes $R_n$, so $ \big( R_n ( f_n)-R_n(f_a)\big)\leq 0$.
\end{frame}

\section{Classification and regression trees}
\frame{\tableofcontents[currentsection]}

\begin{frame}
    \begin{itemize}
	\item $Y_{pred}$ is the random variable representing this prediction (proportional)
	\item $Y$ is the random variable representing the class, in this node (empirical distribution)
    \end{itemize}
    \begin{equation}
        \begin{aligned}
            \label{eq:}
	    P(Y_{pred}\neq Y) &= \sum^{L}_{l=1} P\big(Y_{pred}\neq Y |Y=l\big)P(Y=l)\\
	    &= \sum^{L}_{l=1} \Big(1-P\big(Y_{pred}= Y |Y=l\big)\Big)P(Y=l)\\
	    &= \sum^{L}_{l=1} (1-p_n^l)p_n^l
        \end{aligned}
    \end{equation}
\end{frame}

\begin{frame}{Homogeneity criterion for classification: Gini impurity}
    \begin{equation}
	H(n) = \sum^{L}_{l=1} p_n^l(1-p_n^l)
    \end{equation}
    If we predict the classes in node $n$ according to the proportions of the labels in $n$, then the Gini impurity is the probability of making a mistake, given that we are in node $n$. 
\end{frame}

\end{document} 
